% Dot batch 83; source SHA-256: 6d88758bc1e63aca230d9cdfcd1436d04ecca5bbef1dca8355b61283cfa8a14b
% Archived source: dot_updates_20261124/batches/batch083/EGA第五卷_参数丛单态性与射影标架.tex
% Original body lines: 26--70; integration map: integration/late_integration.json
\begingroup
\def\newscope{V-L}
\def\Spec{\operatorname{Spec}}
\def\Aut{\operatorname{Aut}}
\def\Div{\operatorname{Div}}
\def\OO{\mathcal O}
\def\PP{\mathbf P}
\egas{4.3}{命题}

设 \(D\) 是 \(X/S\) 上由 \(P^\vee\) 参数化的除子线性系，其中 \(g:X\to S\) 平坦且局部有限表示。\cmt{L18-083-01}{法140在“有限表示”上方补写“局部”；(a)的边注另要求 \(g\) 有限表示。英106把两层条件混在一起。本命题的总假设取局部有限表示，(a)再加强，(b)只另要求满射。}

\noindent\textbf{a)}
再设 \(g\) 有限表示。对每个 \(s\in S\)，令 \(k\) 为 \(k(s)\) 的一个代数闭包，且假设几何纤维 \(X_k\) 有一个关联素循环 \(T\)，使自然映射
\[
 k\longrightarrow H^0(T,\OO_T)
\]
是同构。如果 \(g\) 固有且满射，这一条件自动满足。于是在此情形下，态射
\[
 D:P^\vee\longrightarrow\Div_{X/S}
\]
是单态射。\cmt{L18-083-02}{这里 \(T\) 是几何纤维的某个关联点的闭包，随后赋予约化结构；不要求该点是唯一的关联点。固有满射时纤维非空，选取一个不可约分支的泛点；其约化闭包在代数闭域上固有且整，故全局函数为该域，见 \href{https://stacks.math.columbia.edu/tag/04L0}{Stacks，04L0}。法141写 \(\Div_{X/S}\)，英106误写 \(\Div_{S/S}\)。}

\noindent\textbf{b)}
考察映射 \(u\mapsto D\circ u\)，它把 \(\Aut_S(P^\vee)\) 送入 \(X/S\) 上由 \(P^\vee\) 参数化的除子族的集合。如果 \(g\) 满射，这个映射为单射；特别地，\(D=D\circ u\) 蕴含 \(u=\mathrm{id}_{P^\vee}\)。更一般地，自同构函子到以 \(P^\vee\) 参数化的除子线性系函子的态射是单态射。\cmt{L18-083-03}{最后一句据法141手写补入；它指每次换底后的映射也为单射，而不是只比较 \(S\) 的点。用 \(u\) 重参数化仍保持线性系。}

在 a) 的假设下，b) 是 a) 的直接推论；但 b) 只需要较弱的假设。应当注意，若去掉 a) 的额外条件，结论可以不成立。例如取 \(S=\Spec k\)，取 \(\PP^1_k\) 的一个非空开子概形 \(X\)，不含两个不同的 \(k\)-点 \(a,b\)。令 \(D\) 为 \(X\hookrightarrow\PP^1_k\) 所定义的除子族。则 \(a,b\) 在 \(X\) 上定义同一除子，即零除子，但 \(a\ne b\)。\cmt{L18-083-04}{法141明确给出定义除子族的典范浸入；英106省略这一句。此处添明 \(X\ne\varnothing\) 以用非平凡反例。例如取 \(X=\PP^1_k\setminus\{a,b\}\)，这两个点除子分别限制到 \(X\) 均为空，仍是有效 Cartier 除子。}

\medskip\noindent\textbf{域底情形。}
先设 \(S\) 是域 \(k\) 的谱。利用下降，可假定 \(k\) 代数闭。选取 a) 中的 \(T\)，赋予诱导的约化结构。把所得线性系的除子限制到 \(T\)，只须证明相应复合
\[
 P^\vee\longrightarrow\Div_{T/k}
\]
为单态射。限制所得的族仍是除子线性系，因此归结到 \(X=T\)，从而归结到 \(k\xrightarrow{\sim} H^0(X,\OO_X)\) 的情形。\cmt{L18-083-05}{法141和英107都把这一步略写成 \(\Div_{X/k}\to\Div_{T/k}\)；任意除子未必能这样限制，例如在 \(\mathbb A^2_k\) 上把坐标直线作为除子限制到它自身，其方程恒零。因此正文只写这个线性系的限制。\(T\) 的关联泛点处，正则截面不落入那个关联素理想；对几何点同样成立。故在约化整 \(T\) 上得到非零截面，再用相对 Cartier 判据，见本版核查报告 L18-083。}

对任意 \(k\)-概形 \(S'\)，此时有 \(\OO_{S'}\xrightarrow{\sim}g_{S',*}\OO_{X_{S'}}\)。
设 \(L\) 和 \(u:E\to g_*L\) 如4.1。如果 \(E_{S'}\) 的两条处处不为零的截面 \(\phi,\psi\) 所给出的 \(L_{S'}\) 截面 \(u(\phi),u(\psi)\) 有相同的除子，那么两截面相差 \(\OO_{X_{S'}}\) 的一个可逆整体截面。这个单位来自 \(S'\)；于是 \(\phi,\psi\) 也只相差底上的一个单位，从而给出同一个 \(P^\vee(S')\) 的点。每个 \(P^\vee(S')\) 的点在 \(S'\) 局部都由 \(E_{S'}\) 的一条处处不为零的截面给出（见第一章），这就证明了域底情形的 a)。\cmt{L18-083-06}{英107在 \(u\) 的靶处保留问号，按4.1补成 \(g_*L\)。这里用 \(S'\) 区别检验概形和原底 \(\Spec k\)。最后一步还用评估映射对参数截面的单射性；本版核查报告 L18-083 给出有限仿射开检验，覆盖非约化的 \(S'\)。}

\medskip\noindent\textbf{引理。}
设 \(k\) 代数闭，\(X\) 是非空、在 \(k\) 上局部有限型的概形，\(D\) 是 \(X\) 上由 \(P^\vee=\PP(E^\vee)\) 参数化的除子线性系。考虑相应的态射
\[
 f:X-Z\longrightarrow P=\PP(E),
\]
其中 \(Z\) 是基点集。若 \(r=\dim_k E>0\)，则可取 \(r+1\) 个点 \(x_i\in X(k)-Z(k)\)，\(1\leq i\leq r+1\)，使 \(f(x_i)\) 构成 \(P\) 的一个射影标架：任取 \(r\) 个指标，其对应的像不全在一个超平面内。\cmt{L18-083-07}{法142手写补齐 \(x_i\) 和指标范围，英107保留。秩一时 \(P\) 只有一个点、没有射影超平面，标架条件无实质内容，可重用同一点；不可额外要求 \(X\) 在秩一时一定有两个不同的 \(k\)-点。}

可以把 \(X\) 换成 \(X-Z\)，因而假定 \(Z=\varnothing\)。由4.1，\(f(X)\) 不在任何超平面内，因此可先选出 \(r\) 个点 \(x_1,\ldots,x_r\)，其像射影独立，也就是由 \(E\) 上的 \(r\) 个线性无关的线性形式给出。
剩下的是选 \(x_{r+1}\in X(k)\)，使其像避开从这 \(r\) 个像中各取 \(r-1\) 个所张成的 \(r\) 个超平面 \(H_i\)。否则由雅各布森性质可知
\[
 f(X)\subseteq\bigcup_i H_i .
\]
于是，若 \(X_0\) 是 \(X\) 的任一不可约分支，\(f(X_0)\) 必落入某个 \(H_i\)，与4.1矛盾。\cmt{L18-083-08}{法142划去“先把 \(X\) 换成一个整分支”的长段。有效论证是：非零参数截面在各分支的关联泛点均不消失，所以逐分支非退化；仅有 \(f(X)\) 整体非退化，不足以得到射影标架。闭点稠密用 \href{https://stacks.math.columbia.edu/tag/0477}{Stacks，0477}。}

引理既得，为证明域底情形的 b)，也可假定 \(Z=\varnothing\)。由4.1，只须证明知道 \(u\) 在 \(P\) 上的逆转置（逆步）自同构与 \(f:X\to P\) 的复合，就能确定 \(P^\vee\) 的自同构 \(u\)；并且对任意基变换 \(S'\to\Spec k\)，同样能确定 \(P^\vee_{S'}\) 的相对自同构。由刚才的引理和第1小节对 \(\PP(E^\vee)\) 的自同构的描述，这一点随即成立：射影丛的自同构由它在一组射影标架截面上的作用确定，所谓标架截面，是指它们在各纤维上构成射影标架。\cmt{L18-083-09}{法142明确是逆转置（逆步）自同构与 \(f\) 的复合；采用 \(\PP(E)\) 表示一维商的约定。推导时用的是一组在底上定义的截面，不是只知道自同构在每个剩余域上的约化。见本版核查报告 L18-083 的矩阵核查。}
\par\endgroup
